\section{Foundation Technologies}
\label{sec:foundation-technologies}

Stockrium combines financial-analysis methods, large language models, multi-agent coordination, deterministic rules, and historical evaluation. This section briefly introduces the concepts required to understand the system; their implementation is presented in Chapter 4.

\subsection{Financial Analysis Foundations}

Equity analysis commonly considers three complementary sources of evidence: market behavior, company fundamentals, and financial news. Each describes a different aspect of a stock and should therefore be interpreted together rather than treated as an independent guarantee of future performance.

\paragraph{Technical analysis.}

Technical analysis examines historical price and volume data to describe trend, momentum, volatility, and possible turning points \cite{murphy1999technical}. Stockrium uses several complementary indicators. A simple moving average (SMA) is the arithmetic mean of recent closing prices, whereas an exponential moving average (EMA) assigns greater weight to newer observations; both smooth short-term noise and help reveal trend direction. Bollinger Bands place an upper and lower band around a moving average---typically the 20-period SMA plus or minus two standard deviations---so band width reflects recent volatility and the price's position within the bands provides context about relative strength or weakness.

The Relative Strength Index (RSI) is a momentum oscillator bounded between 0 and 100. Values above 70 are conventionally treated as overbought and values below 30 as oversold, although these conditions do not by themselves imply an immediate reversal. Moving Average Convergence Divergence (MACD) uses the difference between the 12-period and 26-period EMAs as its main line, a 9-period EMA of that difference as the signal line, and their difference as the histogram. The relative lines and changing histogram indicate whether momentum is strengthening or weakening. KDJ is a stochastic oscillator named after its K, D, and J lines: K measures the closing price's position within a recent high--low range, D smooths K, and $J=3K-2D$ amplifies the separation between them. Values above 80 or below 20 conventionally indicate overbought or oversold momentum. These indicators are lagging or sensitive to noise, so crossings and threshold values are treated as supporting evidence rather than automatic buy or sell instructions.

\paragraph{Fundamental analysis.}

Fundamental analysis evaluates the financial condition and operating performance of a company using its balance sheet, income statement, cash-flow statement, and derived ratios \cite{penman2013financial}. Stockrium groups non-financial-company indicators into liquidity, leverage, efficiency, profitability, and valuation. The current ratio compares current assets with current liabilities and indicates short-term payment capacity. The debt-to-equity ratio compares liabilities with shareholders' equity and describes financial leverage, while asset turnover measures how efficiently assets generate revenue. Return on Assets (ROA) relates profit to total assets, and Return on Equity (ROE) relates profit to shareholders' equity. Under the DuPont decomposition, ROE can be interpreted as net profit margin multiplied by asset turnover and financial leverage, thereby distinguishing profitability, operating efficiency, and debt-supported returns.

Valuation indicators describe the market price relative to an accounting measure. The price-to-earnings ratio (P/E) compares share price with earnings per share, whereas the price-to-book ratio (P/B) compares market value with book value. A lower ratio is not automatically better because growth expectations, business risk, and sector structure also affect valuation. For financial institutions, Stockrium additionally considers the net interest margin (NIM), capital adequacy ratio (CAR), loan-to-deposit ratio (LDR), and non-performing loan ratio (NPL). These respectively compare net interest income with average earning assets, regulatory capital with risk-weighted assets, gross loans with customer deposits, and non-performing loans with gross loans. Thus, they provide views of interest-earning efficiency, capital resilience, funding liquidity, and asset quality. Compound annual growth rate (CAGR), calculated as $(V_{\mathrm{end}}/V_{\mathrm{start}})^{1/n}-1$ for positive endpoint values separated by $n$ years, summarizes a multi-year trend as a smoothed annual rate. All ratios are interpreted through time trends and industry comparisons because a favorable value in one sector or period may not have the same meaning in another.

\paragraph{Financial news analysis.}

Financial news provides information about corporate events, macroeconomic changes, and market expectations that may not yet appear in accounting data. Sentiment analysis classifies this content as positive, neutral, or negative, while summarization makes long articles easier to review. However, sentiment depends on the affected company, publication time, source, and market context; Vietnamese studies also show that textual polarity alone is not a reliable trading rule \cite{tung2021phobert,vu2023sentiments}. Stockrium therefore treats news as supporting evidence alongside technical and fundamental information.

\subsection{Large Language Models}

Large language models (LLMs) are neural models trained on extensive text collections to understand and generate natural language. Most modern LLMs are based on the Transformer architecture, which uses attention to represent relationships among words in a context \cite{vaswani2017attention}. Their instruction-following capability allows one model to summarize articles, interpret financial observations, classify content, and generate explanations without training a separate model for every task \cite{brown2020gpt3}. Nevertheless, LLMs can produce unsupported or inconsistent statements, and a valid response structure does not guarantee factual or numerical correctness. In Stockrium, LLMs are consequently used mainly for interpretation and explanation, while exact calculations and selected validation rules remain in conventional code.

\subsection{AI Agents, Multi-Agent Systems, and Coordination Workflows}

An artificial-intelligence (AI) agent combines a model with a defined role, available information, state, and actions such as calling an application programming interface (API) or returning a structured result. The Reasoning and Acting (ReAct) pattern illustrates how model reasoning can be interleaved with external actions and their observations \cite{yao2022react}. A multi-agent system divides a larger task among cooperating agents \cite{wooldridge2009multiagent}. This separation allows each Stockrium agent to focus on technical data, fundamentals, or news and exposes intermediate results for later synthesis. The approach improves modularity but can also increase latency and propagate errors, making clear responsibilities and compatible output formats necessary.

Several common patterns organize agent-based applications. A \textit{specialized agent} owns one domain; an \textit{orchestrator--worker} structure assigns work to the appropriate agents; \textit{fan-out/fan-in} executes independent branches in parallel before joining them; and an \textit{aggregator} reconciles their outputs. Shared state provides a common data contract, while conditional routing selects the next node from the current state. Stockrium combines these patterns to make its analytical flow explicit and bounded rather than relying on unrestricted agent-to-agent conversation.

A multi-agent workflow represents processing as a graph of nodes and transitions. Nodes may retrieve data, perform calculations, invoke an LLM, or update shared state, while edges define sequential, conditional, or parallel execution. LangGraph provides these graph and state abstractions together with checkpointing for persistent conversations \cite{langgraph2026graphapi}. In Stockrium, one graph coordinates the stock-analysis agents and another routes chatbot requests. This structure makes the allowed execution paths visible and allows deterministic functions and language-model nodes to operate within the same workflow.

\subsection{Hybrid LLM--Rule Decision Support}

A hybrid decision-support system assigns qualitative interpretation to LLMs and reproducible calculations to deterministic components. Stockrium's agents interpret evidence and propose structured assessments, whereas conventional code calculates indicators and a normalized confidence score, checks minimum conditions, and constrains selected recommendation values according to the investment horizon. The confidence score lies between 0 and 1 and is a weighted aggregation of the active technical, fundamental, and news components; it is an internal strength-of-evidence measure, not a statistically calibrated probability of profit. This arrangement improves consistency and traceability but does not eliminate uncertainty because the inputs and qualitative interpretations may still be incomplete or incorrect. The generated result must therefore be understood as decision support rather than a guaranteed investment outcome.

\subsection{Backtesting and Evaluation Foundations}

Backtesting applies a defined rule set or analytical pipeline to historical data to estimate how it would have behaved under past market conditions. A meaningful simulation specifies signal, entry, exit, holding-period, and transaction-cost assumptions while preventing the use of information that was unavailable at the simulated decision time. In this thesis, total return compounds successive closed-trade returns as $\prod_{i=1}^{N}(1+r_i)-1$ for each security; win rate is the proportion of the $N$ closed trades with a positive net return; and maximum drawdown is the largest peak-to-trough decline in the simulated equity curve. Trade count is reported because percentages derived from very few observations are less reliable.

The Sharpe ratio generally measures excess return per unit of return variability \cite{sharpe1994ratio}. Stockrium's implementation uses closed-trade returns, a zero risk-free benchmark, and computes $\sqrt{N}\,\bar{r}/s_r$, where $\bar{r}$ and $s_r$ are the sample mean and sample standard deviation of trade returns. It returns zero when fewer than two trades are available or when their returns have no variation. The result is therefore a trade-level comparative statistic rather than a conventional annualized daily-return Sharpe ratio. Simple baselines provide reference behavior, market-regime analysis separates results under different broad market trends, and walk-forward evaluation repeats testing over successive time windows instead of relying on one fixed split. Nevertheless, historical results can be distorted by overfitting, survivorship bias, look-ahead bias, incomplete data, and simplified execution assumptions \cite{bailey2017pbo}; they describe past behavior and do not prove future profitability. Stockrium's backtesting facility is consequently an internal evaluation tool for its pipeline and supported configurations, not a general-purpose strategy-construction platform.

Together, these foundations explain the principal division of responsibilities in Stockrium: financial methods define the evidence, LLMs interpret it, the multi-agent system coordinates specialized work, deterministic rules constrain the result, and backtesting evaluates historical behavior.
